\Needspace{8\baselineskip}\originalchapter{2002 年期末卷：动力学、束与平衡}
2002 Fall 22.611J 期末原题1—6 每组有官方扫描答案。以下将题面主条件完整写出、解法逐项给出；对原题未给出的束密度和碰撞参数保留符号，不能据此虚构唯一增长率。
\begin{exercise}\label{ex:hist02-final-1}
原题1（18分）：(a—c)写电子／$Z_i$价离子的等离子体频率、离子回旋频率、电子／离子 Debye 长度；(d—i)判断 Vlasov 方程是否表示六维相空间不可压缩、相空间粒子数守恒、三维流体速度 $\mathbf V=n^{-1}\int\mathbf vf\dd^3v$ 不可压缩、精细熵守恒、时间可逆、均匀分布自动趋于 Maxwell。
\end{exercise}
\begin{solution}
编者独立解答。$\omega_{pe}^2=n_ee^2/(\epsilon_0m_e)$、$\omega_{pi}^2=n_iZ_i^2e^2/(\epsilon_0m_i)$、$\Omega_i=Z_ieB/m_i$、$\lambda_{Ds}^2=\epsilon_0T_s/(n_sq_s^2)$。Vlasov 特征流 $\dot{\mathbf x}=\mathbf v$、$\dot{\mathbf v}=(q/m)(\mathbf E+\mathbf v\times\mathbf B)$ 的六维散度为零，且 $Df/Dt=0$；故 d、e 真。三维 $\nabla_x\cdot\mathbf V$ 可非零，f 假。无边界通量时 $\dd\int f\ln f\dd^3x\dd^3v/\dd t=0$，g 真（指精细熵），把速度、电磁时间反演条件一同反向则方程可逆，h 真。无碰撞的均匀 $f(\mathbf v)$ 是定态，不会自发 Maxwell 化，i 假。单独局域位置的 $-\int f\ln f\dd^3v$ 可随输运改变，须区分与全域熵。
\end{solution}

\begin{exercise}\label{ex:hist02-final-2}
原题2（10分）：仅知实频上的解析介电函数 $\epsilon(k,\omega_r)$，在 $|\gamma|\ll\omega_r$ 时求 $\omega=\omega_r+\ii\gamma$ 的实频和阻尼／增长，解释波能与耗散，并将公式用于原卷下题式(1)的纵向介电函数 $\epsilon=1-(\omega_{pe}^2/k^2)\int_L F'(u)/(u-\omega/k)\dd u$。
\end{exercise}
\begin{solution}
编者独立解答。先解 $\operatorname{Re}\epsilon(k,\omega_r)=0$。在 $\omega_r$ 展开 $\epsilon(k,\omega_r+\ii\gamma)\simeq\operatorname{Re}\epsilon+\ii[\operatorname{Im}\epsilon+\gamma\partial_\omega\operatorname{Re}\epsilon]$，从而
\begin{equation}\gamma=-\frac{\operatorname{Im}\epsilon(k,\omega_r)}{\partial_\omega\operatorname{Re}\epsilon(k,\omega_r)}.
\end{equation}
电纵波时均能量 $W=(\epsilon_0|E|^2/4)\partial_\omega(\omega\operatorname{Re}\epsilon)$，吸收功率 $Q=(\epsilon_0\omega_r|E|^2/2)\operatorname{Im}\epsilon$；在根上 $2\gamma W=-Q$，负波能时“正耗散”也可增幅。Landau 处方给式(1) $\operatorname{Im}\epsilon=-\pi(\omega_{pe}^2/k^2)F'(\omega_r/k)$ 或相反符号，取决于题卷 $L$ 绕极点的方向；本卷采用 $e^{-\ii\omega t}$ 与 $\gamma>0$ 增长的约定时，正斜率分布对应 $\gamma>0$。必须同时核对 $F'$ 与实部导数符号，不能孤立地把虚部称为阻尼率。
\end{solution}

\begin{exercise}\label{ex:hist02-final-3}
原题3（32分）：背景电子 $T_e\simeq1\,$keV 中注入两条稀薄电子束，中心能 $E_{b1}=10\,$keV、$E_{b2}=100\,$keV，束温宽 $T_b\simeq1\,$keV，离子冷且动态忽略。定义背景热速 $v_e$、束宽 $v_b$、束中心 $U_{b1},U_{b2}$，正 $k$；用峰左端 $v_{m,j}=U_{bj}-\alpha v_b$（$\alpha\sim1$）辨增长带，求领先阶最不稳定根的波能和增长率，比较长期线性主导模。
\end{exercise}
\begin{solution}
编者独立解答。按 $f\propto e^{-(v-U)^2/v_t^2}$ 的题卷热速，$v_e=v_b=\sqrt{2(1\mathrm{keV})/m_e}\simeq1.88\times10^7\,$m/s，非相对论 $U_{b1}=\sqrt{10}v_b$、$U_{b2}=10v_b$；后者 $\simeq0.63c$，定量预测需相对论修正。两束左坡 $v_m<\omega_r/k<U_b$ 可反 Landau 阻尼；背景在 $\omega_r/k\sim U_{b1}$ 仍有尾部阻尼，因此若束密度任意趋零，第一带不保证真的增长。对每条束在 $n_{bj}\ll n_p$ 且其驱动超过背景阻尼时，$\omega_r\simeq\omega_{pe}$、$k_j\simeq\omega_{pe}/(U_{bj}-v_b/\sqrt2)$；以 $\xi_j=(\omega_r/k-U_{bj})/v_b$，
\begin{equation}\gamma_j\simeq-\sqrt\pi\,\frac{n_{bj}}{n_p}\frac{\omega_{pe}^3}{k_j^2v_b^2}\xi_j e^{-\xi_j^2}+\gamma_{p,j},\quad \xi_j\simeq-1/\sqrt2,\quad \gamma_{p,j}\le0.
\end{equation}
近 Langmuir 根 $\partial_\omega(\omega\epsilon_R)\simeq2$，故 $W_j\simeq\epsilon_0|E_j|^2/2$、按 $\epsilon_0|E|^2/4$ 归一的无量纲波能为2。若两束密度相近且都不受背景阻尼主导，第二带 $\gamma_2/\gamma_1\sim(U_{b2}/U_{b1})^2\sim10$，长期线性先占优；真实长期饱和要考虑准线性展平，不能由线性式判定。原题没给 $n_{b1},n_{b2}$，所以只能给比例和条件，不能给唯一 s$^{-1}$ 数值。
\end{solution}

\begin{exercise}\label{ex:hist02-final-4}
原题4（15分）：从电子与离子的 Maxwell--Vlasov 介电函数 $\epsilon=1+[1+\zeta_eZ(\zeta_e)]/(k^2\lambda_{De}^2)+[1+\zeta_iZ(\zeta_i)]/(k^2\lambda_{Di}^2)$ 出发，其中 $v_s=\sqrt{2T_s/m_s}$、$\zeta_s=\omega/(kv_s)$；在 $v_i\ll\omega/k\ll v_e$ 与 $k\lambda_{De}\ll1$ 下求离子声色散并证明准中性。
\end{exercise}
\begin{solution}
编者独立解答。$|\zeta_e|\ll1$ 使电子实响应近 $1/(k^2\lambda_{De}^2)$；$|\zeta_i|\gg1$ 时 $1+\zeta_iZ(\zeta_i)\simeq-1/(2\zeta_i^2)$，故 $\epsilon_R\simeq1+1/(k^2\lambda_{De}^2)-\omega_{pi}^2/\omega^2$。解得 $\omega^2=k^2(T_e/m_i)/(1+k^2\lambda_{De}^2)$，长波 $\omega\simeq kc_s$、$c_s^2=T_e/m_i$。Poisson 式 $e(\delta n_i-\delta n_e)=\epsilon_0k^2\phi$，电子 $\delta n_e=n_0e\phi/T_e$，故 $(\delta n_i-\delta n_e)/\delta n_e=k^2\lambda_{De}^2\ll1$。$T_i\ll T_e$ 还使离子 Landau 阻尼指数小，才有清晰的声波根。
\end{solution}

\begin{exercise}\label{ex:hist02-final-5}
原题5（10分）：原卷已给 Maxwell 电子在离子背景下的能量交换矩 $\int(m_ev^2/2)C_{ei}^{E}\dd^3v=-\nu_{ei}(m_e/m_i)\sqrt{2/\pi}\,n(T_e-T_i)$。不重做该积分，只用碰撞算子的基本守恒性质求 $C_{ie}$ 的离子能量矩并解释能量流向。
\end{exercise}
\begin{solution}
编者独立解答。交叉碰撞逐次守恒总能量，且种内算子各自能量矩为零。因此
\begin{equation}\int\frac{m_iv^2}{2}C_{ie}\dd^3v=+\nu_{ei}\frac{m_e}{m_i}\sqrt{\frac2\pi}\,n(T_e-T_i).
\end{equation}
$T_e>T_i$ 时电子失去、离子获得同量能；反向温差时符号反转。该传热比电子角散射慢约质量比，不能从“总能量守恒”推出每个物种各自能量不变。题卷的 $\nu_{ei}$ 和热速规范按原题式(2)使用，避免与别章不同热速约定误比前因子。
\end{solution}

\begin{exercise}\label{ex:hist02-final-6}
原题6（15分）：(a) 在固定 $n,E$ 下极大化 $S=-\int f\ln f\dd^3v$ 求 Maxwell；(b) 同种 Landau 碰撞算子作用时判 $S$ 变化方向；(c) 对原题所给径向电子—离子能量算子 $C_{ei}^{E}$ 求唯一正则静态分布，并解释拖曳和扩散的平衡。
\end{exercise}
\begin{solution}
编者独立解答。(a) 令 $\delta[S+\alpha n+\beta E]=0$ 得 $\ln f=\alpha-1+\beta mv^2/2$，归一与能量约束给 $f=n[m/(2\pi T)]^{3/2}e^{-mv^2/(2T)}$；$\delta^2S=-\int(\delta f)^2/f<0$。(b) 把 $\dot S=-\int C(f,f)\ln f$ 与交换 $\mathbf v,\mathbf v'$ 后的式子平均，可化为
\begin{equation}\dot S=\frac\Gamma2\iint ff'\,[\nabla_v\ln f-\nabla_{v'}\ln f']\cdot U(\mathbf v-\mathbf v')\,[\nabla_v\ln f-\nabla_{v'}\ln f']\dd^3v\dd^3v'\ge0,
\end{equation}
因为 $U$ 半正定。(c) 无源、可积且边界通量为零，径向速度空间流必须为零；题卷式(2)内括号令 $\partial_v f=-(m_ev/T_i)f$，解 $f\propto e^{-m_ev^2/(2T_i)}$。拖曳把电子速度拉低，速度扩散由热离子热浴注入涨落，两项恰在 $T_e=T_i$ 互相抵消；若仅保留拖曳，不会得到 Maxwell 稳态。
\end{solution}
